\begin{abstract}
Early transformer layers map document tokens into intermediate semantic representations that later layers can use for prediction. Our key insight is to reuse this encoder-like computation across queries by placing a persistent cache between the two parts of the decoder. CoMem (Comprehension Memory) treats the lower layers as a semantic encoder and the remaining layers as a reader. Write encodes document chunks once and stores their intermediate residuals; Read retrieves selected cached states, combines them with the encoded query, and continues through the upper layers. A self-distilled suffix adapter learns to consume independently encoded states. On Qwen3-8B, this cache uses $1/18$ of the payload storage of full-depth KV. With identical evidence and the same adapter, CoMem provides a $1.403\times$ selected-pack prefill speedup at a $3.12$-point RULER accuracy cost. The semantic-encoder view thus motivates an intermediate cache that saves repeated document computation; its end-to-end benefit depends on reuse frequency, preparation cost, and task-specific representation fidelity.
\end{abstract}
